% Terminal note of the homily: everything the speech itself must not carry ---
% the audience, the spoken word count and the pace it implies, the relation to
% the three interpretations of the expansive study, the route by which each
% quoted or reported text reaches the page, and the exact loci --- followed by
% the References. Nothing here is read aloud.

\section*{Note on Sources and Delivery}

{\small

\textbf{Audience and occasion.} An adult parish assembly, at the Mass of the
Twenty-ninth Sunday in Ordinary Time, Year~A: \latin{Dominica XXIX ``per
annum''} of the \work{Roman Missal}, Third Edition, for the dioceses of the
United States, with no.~145 of the \work{Lectionary for Mass}, for the
occurrence of 18~October 2026. Every appointed element is required except the
Communion antiphon, of which the Missal gives two; the speech names the second
only as one of the two, through the verse of Mark it takes up, so that what it
says holds whichever is sung.

\textbf{Length and pace.} The spoken body runs to 1,466 words, counted as
whitespace-separated words once comments and markup are removed: a little over
eleven minutes at an unhurried 130 words a minute, a little under twelve at
125, and a little over twelve at 120. That is arithmetic on the word count; no
speaker has delivered these words and no delivery has been timed. The prose
was read through in full in its final form for sense, sentence length and ease
of speech, silently and not aloud.

\textbf{Relation to the reviewed interpretations.} The argument preached joins
the first and third of the expansive study's three interpretations, which its
comparison finds can be preached together without strain. From the first,
\textit{Render to God the One Who Bears His Image}, it takes the inference
from the coin's image to the human person made in God's image, with
Augustine's ``to God yourselves'', the coin worn smooth and refashioned, and
Ambrose's Christ who bears no image of Caesar; Jerome's literal answer, tithes
and offerings, is kept beside it as a fair reading of the words. From the
third, \textit{True Words from a Divided Heart}, it takes the hypocrite whose
true words come from a divided heart, with Chrysostom and Aquinas, Chrysostom's
proof of Godhead from the concealed intent, and the Collect's sincere heart
with the acclamation's sentence. From the second, \textit{The Lord Who Grasps
the Hand of Kings}, it takes only what all three share, that Caesar's claim
stands beneath God's: the Lord's hand on Cyrus, Chrysostom's render and not
give and his limit at godliness, and Augustine's soldiers under Julian. The
witnesses' disagreement over the ground of Caesar's claim is not preached; the
speech says only that the questioners were carrying Caesar's coin, which
Hilary and Ambrose, who ground the debt in possession, and Chrysostom and
Augustine alike accept. The christological reading of Christ's knowledge of
their hearts is credited to Chrysostom alone. The agents of the scene are
those the Gospel names, the vice is applied to the assembly, and none of the
polemical sentences that the research records exclude under \work{Nostra
aetate}~4 is used.

Each Father is credited only with what he says of his own passage:
Augustine's double noose; Chrysostom's question asked so that they would be
bound by their own answers, his render and not give, his ``not with an upright
mind'', his rebuke that was no conjecture but a sign of knowing their secret
thoughts and a proof of Godhead, given with gentleness, his limit at
godliness, and his labour of love; Jerome's tithes, first-fruits and offerings;
Augustine's God who asks for his money as Caesar for his, his coin worn smooth
by wandering and refashioned, his stamping by a new birth, his soldiers under
Julian and his humble heart as the present for God's courts; Ambrose's Christ
who bears no image of Caesar, his renunciation in words and not in affection
when the sacraments are received, his contract of faith greater than that of
money and his ``render what you promised''; Aquinas's hypocrite with one thing
in the mouth and another in the heart, and his single intention of the
children of a simple God. The following connections are the homily's own, or
the studies' editorial synthesis, and no Father makes them: Ambrose's Christ
as the one who came to strike Augustine's worn coin again, and his self-gift
to the Father with Mark 10:45; the restamping placed at the font with
Ambrose's promise; Ambrose's renunciation in words and not in affection set
beside the questioners' true words; the observation that Ambrose's
\latin{Redde} is the Gospel's \latin{Reddite}; the Collect and the
acclamation's \latin{simplices} as the contrary of hypocrisy; Cyrus beside
Caesar, so that Caesar's coin is minted in God's world, which is the second
interpretation's inference from the Lectionary's pairing; the bread and wine
as God's gifts with the offerers going up with them; the Prayer after
Communion's temporal and eternal goods as the coin that helps and does not
teach; Communion as the restamping of the image; and the practicable response.

Left to the studies: Hilary's body, soul and will and his ground of the debt;
Aquinas's literal and mystical senses, his answer first to the mind, and his
``You are God's, and Caesar's''; Jerome's Christ who rendered by doing the
Father's will, and his rule of discernment for any answerer; Cyril; the
counter-question as the witnesses' convergence; the ``Christ my Lord'' variant
and the typology of \latin{assimilavi}; the psalm's gods of the nations and
Augustine's ``from the wood''; the Entrance Antiphon; the first Communion
antiphon; the Thessalonian church under Caesar's decrees; the stars of
Philippians; Justin and Romans~13; Bellarmine and Bede; and the sacramentary
history of the orations.

\textbf{Texts quoted, and their route.} Scripture is quoted in the
Douay--Rheims version as revised by Bishop Challoner, a public-domain study
translation, at Mt 22:16--21, Gn 1:26, Is 45:1, 4--6, Ps 95:8, Phil 2:15,
Mk 10:45 and 1~Jn 3:2. It is not the English proclaimed at Mass, which is that
of the \work{Lectionary for Mass for Use in the Dioceses of the United States
of America} and is not reproduced; the speech says once that the verb render
is the old translation's. The Collect, the Prayer over the Offerings and the
Prayer after Communion are described after the study's descriptions of the
Latin of the \work{Missale Romanum}; the descriptions are not translations and
are not the prayers' words. Augustine's tractates and expositions of the
psalms and Chrysostom's homilies on Matthew and on First Thessalonians are
quoted in the English of the Nicene and Post-Nicene Fathers; each quotation
was read in the tracked text at the locus below. Augustine's double noose and
his new birth are reported in that English without quotation marks. Ambrose,
Jerome and Aquinas are reported without quotation marks, in English that
renders the Latin the study prints; Ambrose's Latin was read again in the
tracked electronic text.

\textbf{Exact loci.} Augustine, the double noose: \work{Enarrationes in
Psalmos} 57.8; man stamped again ``by a new birth'': 4.8; ``Enter with an
humble heart into the house of God, and thou hast entered with a present'':
95.9; the soldiers under Julian and ``They distinguished their everlasting
from their temporal master'': 124.7; ``We are God's money~\ldots'', ``He has
come to refashion~\ldots'', God ``asking for His money, as Caesar for his'' and
``to Caesar his money, to God yourselves'': \work{In Iohannis evangelium
tractatus} 40.9, where Augustine also cites 1~John 3:2. Chrysostom, the praise
``not with an upright mind'' and the rebuke as ``a sign of His knowing their
secret thoughts'': \work{Homilies on Matthew} 70.1; ``bound by their own
answers'', ``this is not to give but to render'', ``no longer Caesar's tribute,
but the devil's'' and the ``proof of His Godhead'': 70.2; ``Merely to love is
no labor at all~\ldots'': \work{Homilies on First Thessalonians}~1. Jerome,
\latin{decimas primitias, et oblationes, ac victimas}: \work{Commentarii in
Matthaeum} III, on Mt 22:21, PL 26, col.~163. Ambrose, \latin{Imaginem
Caesaris non habet Christus, quia imago est Dei}: \work{Expositio evangelii
secundum Lucam} IX.35; \latin{omnes renuntiamus verbis, sed non omnes
renuntiamus affectu; nam cum sacramenta suscipimus, renuntiamus \ldots\ Maior
est contractus fidei, quam pecuniae. Redde promissum, dum in hoc corpore es}:
IX.36. Aquinas, \latin{hypocritae sunt proprie qui aliud in ore habent, aliud
in corde}: \work{Super Matthaeum} c.~22, Venice ed., p.~283; \latin{Deus
autem simplex est: unde simplices simus sicut filii Dei; quod est, quando
intentio est ad unum}: \work{Super ad Philippenses} c.~2, lect.~4, Dessain
ed., p.~383.

\par}

% The References stand whole on one page where they can: a list broken after
% one or two entries reads as though the rest had been lost.
\Needspace{20\baselineskip}
\section*{References}

{\small
\begin{itemize}\setlength{\itemsep}{0.15em}\setlength{\parskip}{0pt}
\item \work{Missale Romanum}, editio typica tertia (Typis Vaticanis, 2002), \latin{Dominica XXIX ``per annum''}.
\item \work{Ordo lectionum Missae}, editio typica altera (Libreria Editrice Vaticana, 1981), no.~145, and \work{Lectionary for Mass for Use in the Dioceses of the United States of America}, second typical edition, vol.~I, no.~145, as witnessed by the USCCB \work{Liturgical Calendar} for 2026 and page of readings for 18~October 2026.
\item The Holy Bible, Douay--Rheims version, revised by Bishop Richard Challoner (Project Gutenberg text): Gn 1:26; Ps 95:8; Is 45:1, 4--6; Mt 22:16--21; Mk 10:45; Phil 2:15; 1~Jn 3:2.
\item \work{Biblia Sacra Vulgatae editionis} (Clementine Vulgate), eBible.org text: Mt 22:17, 21; Phil 2:15.
\item St Augustine, \work{Enarrationes in Psalmos} 4.8; 57.8; 95.9; 124.7; Nicene and Post-Nicene Fathers, 1st ser., vol.~8.
\item St Augustine, \work{In Iohannis evangelium tractatus} 40.9, Nicene and Post-Nicene Fathers, 1st ser., vol.~7.
\item St Ambrose, \work{Expositio evangelii secundum Lucam} IX.35--36 (Migne, PL 15), in an electronic Latin text of that printing.
\item St Jerome, \work{Commentarii in Matthaeum} III, on 22:21 (Migne, PL 26, Paris, 1845, col.~163).
\item St John Chrysostom, \work{Homilies on Matthew} 70.1--2, Nicene and Post-Nicene Fathers, 1st ser., vol.~10; \work{Homilies on First Thessalonians} 1, vol.~13.
\item St Thomas Aquinas, \work{Super Evangelium Matthaei} c.~22, in \work{Opera}, editio altera Veneta, vol.~3 (Venice, 1745), p.~283; \work{Super epistolam ad Philippenses} c.~2, lect.~4, in \work{In omnes S.~Pauli Apostoli epistolas commentaria}, vol.~2 (Li\`ege: Dessain, 1857), p.~383.
\end{itemize}
\par}
